
\medskip
Let \((\mathrm{Ab})\)\footnote{\label{sga3:VIIB:nde:10}The editors
inserted \SGARefLink{sga3:VIIB:0.2.E}{Proposition 0.2.E} and \SGARefLink{sga3:VIIB:0.2.F}{Corollary 0.2.F} into this paragraph;
they will be useful in \SGARefLink{sga3:VIIB:0.2.2}{0.2.2}. In the original, these results appeared
in \SGARefLink{sga3:VIIB:0.3}{0.3}.} be the category of abelian groups and let
\(\operatorname{LF}(A)\) be the full subcategory of
\(\operatorname{PC}(A)\) formed by the objects of finite length. For
every object \(M\) of \(\operatorname{PC}(A)\), let \(h_c^M\) denote
the functor
\[
  \operatorname{LF}(A)\longrightarrow(\mathrm{Ab}),
  \qquad
  N\longmapsto\operatorname{Hom}_c(M,N).
\]
By \textup{[CA]}, \S~II.4, Theorem~1, Lemma~4, and Corollary~1, we have
the following results.\footnote{\label{sga3:VIIB:nde:11}\textit{Editors' note.} Indeed,
\(\operatorname{PC}(A)\) has a set of Artinian cogenerators, filtered
projective limits are exact in it, and \(\operatorname{LF}(A)\) is the
subcategory of Artinian objects. The dual category
\(\operatorname{PC}(A)^\circ\) therefore has a set of Noetherian
generators, and filtered inductive limits are exact in it. By the proof
of \textup{[CA]}, \S~II.4, Theorem~1, the functor
\(M\mapsto\operatorname{Hom}_c(M,-)\) is an anti-equivalence from
\(\operatorname{PC}(A)\) to
\(\operatorname{Lex}(\operatorname{LF}(A),(\mathrm{Ab}))\).
Likewise, Lemma~4 and Corollary~1 of \textit{loc.\ cit.} state results
``dual'' to those of \SGARefLink{sga3:VIIB:0.2.F}{Corollary~0.2.F}. For a ``direct'' proof, see
\textup{[DG70]}, \S~V.2, Theorem~3.1, Lemma~3.5, and
Corollaries~3.3 and~3.4.}

\medskip
\noindent\textbf{Proposition 0.2.E.}
\phantomsection
\label{sga3:VIIB:0.2.E}
\textit{The functor \(M\mapsto h_c^M\) is an anti-equivalence from
\(\operatorname{PC}(A)\) to the category}
\[
  \operatorname{Lex}\bigl(\operatorname{LF}(A),(\mathrm{Ab})\bigr)
\]
\textit{of left-exact functors
\(\operatorname{LF}(A)\longrightarrow(\mathrm{Ab})\).}

\medskip
\noindent\textbf{Corollary 0.2.F.}
\phantomsection
\label{sga3:VIIB:0.2.F}
{\itshape
\begin{enumerate}
  \item[\textup{(i)}]
  An object \(P\) of \(\operatorname{PC}(A)\) is projective if and
  only if the functor \(h_c^P\) is exact; that is, if and only if the
  functor \(\operatorname{Hom}_c(P,-)\) is exact on
  \(\operatorname{LF}(A)\).

  \item[\textup{\SGAInlineRefTarget{sga3:VIIB:0.2.F:item:ii}(ii)}]
  Let \((M_i)\) be a filtered projective
  system of objects of \(\operatorname{PC}(A)\).
  For every object \(N\) of \(\operatorname{LF}(A)\), there is an
  isomorphism, functorial in \(N\),
  \[
    \operatorname{Hom}_c\!\left(\varprojlim_i M_i,N\right)
    \simeq
    \varinjlim_i\operatorname{Hom}_c(M_i,N).
  \]

  \item[\textup{\SGAInlineRefTarget{sga3:VIIB:0.2.F:item:iii}(iii)}]
  Every filtered projective limit and every
  product of projective objects of
  \(\operatorname{PC}(A)\) is a projective object of
  \(\operatorname{PC}(A)\).
\end{enumerate}}

Finally, \SGARefLink{sga3:VIIB:0.2.F}{Corollary~0.2.F} gives the following result.

\medskip
\noindent\textbf{Corollary 0.2.G.}
\phantomsection
\label{sga3:VIIB:0.2.G}
\textit{Let \((M_i)_{i\in I}\) be a family of objects of
\(\operatorname{PC}(A)\). Then \(\prod_{i\in I}M_i\) is projective if
and only if each \(M_i\) is projective.}

Indeed, for every \(N\in\operatorname{Ob}\operatorname{LF}(A)\), one
has
\[
  \operatorname{Hom}_c\!\left(\prod_i M_i,N\right)
  \simeq
  \bigoplus_i\operatorname{Hom}_c(M_i,N).
\]

\medskip
\noindent\textbf{0.2.1.}
\phantomsection
\label{sga3:VIIB:0.2.1}
Each local component \(A_{\mathfrak m}\) of \(A\) is a direct factor
of \(A\), hence a projective object of \(\operatorname{PC}(A)\)
(\(A\) is manifestly projective). Moreover,
\[
  S_{\mathfrak m}
  =
  A_{\mathfrak m}/\mathfrak mA_{\mathfrak m}
\]
is the unique simple quotient of \(A_{\mathfrak m}\), so
\(A_{\mathfrak m}\) is indecomposable. On the other hand, every simple
object of \(\operatorname{PC}(A)\) is isomorphic to a unique
\(S_{\mathfrak m}\). By \textup{[CA]}, IV, \S~3, Corollary~1 of
Theorem~3, we
therefore have:

\begin{samepage}
\medskip
\noindent\textbf{Proposition.}
\phantomsection
\label{sga3:VIIB:0.2.1:projectives}
{\itshape
\begin{enumerate}
  \item[\textup{(i)}]
  Every projective object of \(\operatorname{PC}(A)\) is a direct
  product of indecomposable projective objects, uniquely determined up
  to isomorphism.

  \item[\textup{\SGAInlineRefTarget{sga3:VIIB:0.2.1:item:ii}(ii)}]
  Every indecomposable projective object is isomorphic to a unique
  \(A_{\mathfrak m}\), where
  \(\mathfrak m\in\Upsilon(A)\).
\end{enumerate}}
\end{samepage}

\medskip
\noindent\textbf{Definition.}
\phantomsection
\label{sga3:VIIB:0.2.1:topologically-free}
A pseudocompact \(A\)-module \(M\) is said to be
\emph{topologically free} if it is isomorphic to the product of a
family \((A_i)\) of copies of \(A\).

In this case, a family \((m_i)\) of elements of \(M\) is called a
\emph{pseudobasis} of \(M\) if the \(A\)-linear maps
\(A_i\longrightarrow M\) sending the unit element of \(A_i\) to
\(m_i\) extend to an isomorphism
\[
  \prod_i A_i\xrightarrow{\ \sim\ }M.
\]

\medskip
\noindent\textbf{0.2.2.}
\phantomsection
\label{sga3:VIIB:0.2.2}
\footnote{\label{sga3:VIIB:nde:14}The editors have spelled out the
results of this paragraph, which play an important role below; cf.\
\SGARefLink{sga3:VIIB:1.2.3}{1.2.3}, \SGARefLink{sga3:VIIB:1.3.5}{1.3.5}, \SGARefLink{sga3:VIIB:2.2.1}{2.2.1}, and so on.}
If \(M\) is a pseudocompact \(A\)-module, we write \(M^\dagger\) for
the \(A\)-module, without topology,
\[
  \operatorname{Hom}_c(M,A).
\]

Conversely, if \(N\) is an \(A\)-module, we write
\[
  N^*=\operatorname{Hom}_A(N,A)
\]
for its dual, endowed with the topology of pointwise convergence. Thus
a basis of neighborhoods of \(0\) in \(N^*\) is formed by the
submodules
\[
  \mathcal V(x,\mathfrak l)
  =
  \{f\in N^*\mid f(x)\in\mathfrak l\},
\]
where \(x\in N\) and \(\mathfrak l\) is an open ideal of \(A\).

This makes \(N^*\) a pseudocompact \(A\)-module. Indeed, first suppose
that \(N=A\). Then \(N^*=A\), endowed with its topology as a
pseudocompact ring. If \(N\) is a free module \(A^{(I)}\), then
\(N^*\) is the product \(A^I\), endowed with the product topology.
Furthermore, for every morphism
\(\varphi:N_1\longrightarrow N_2\), the transpose
\[
  {}^t\!\varphi:N_2^*\longrightarrow N_1^*
\]
is continuous, because the inverse image under \({}^t\!\varphi\) of
a submodule \(\mathcal V(x_1,\mathfrak l)\) of \(N_1^*\) is precisely
the submodule \(\mathcal V(\varphi(x_1),\mathfrak l)\) of \(N_2^*\).
For arbitrary \(N\), choose a presentation
\[
  A^{(J)}
  \xrightarrow{\ \varphi\ }
  A^{(I)}
  \xrightarrow{\ \pi\ }
  N
  \longrightarrow
  0.
\]
Then \(N^*\) is the kernel of the continuous morphism
\[
  {}^t\!\varphi:A^I\longrightarrow A^J,
\]
and hence \(N^*\) is pseudocompact.

When \(A\) is Artinian (in which case one may take
\(\mathfrak l=0\) above), \SGARefLink{sga3:VIIB:0.2.F}{Corollary~0.2.F} gives:

\medskip
\noindent\textbf{Proposition.}
\phantomsection
\label{sga3:VIIB:0.2.2:artinian-duality}
\textit{When \(A\) is Artinian, the functors}
\[
  P\longmapsto P^\dagger
  \qquad\textit{and}\qquad
  Q\longmapsto Q^*,
\]
\textit{where \(P\), respectively \(Q\), is a projective object of
\(\operatorname{PC}(A)\), respectively a projective \(A\)-module,
establish an anti-equivalence between the category of projective
pseudocompact \(A\)-modules and the category of projective
\(A\)-modules.}

In particular, when \(A\) is a field \(k\), the functor
\[
  P\longmapsto P^\dagger
\]
is an anti-equivalence between the category of all pseudocompact
\(k\)-modules (also called \emph{linearly compact \(k\)-vector
spaces}) and the category of \(k\)-vector spaces.

